% Section 6.4, from lectures/L06/appendix/b_exercises.md (Set 1) and c_conv_layer.md (Set 2). The
% hand-training example's walk-through is the worked solution in Appendix A; the conv layer is
% finished in Chapter 7, and its reference implementations are in the repository.

\section{Exercises}
\label{c6:sec:exercises}

\begin{aside}
\textbf{How to work these.} Set~1 is a calculation: work it on paper, Exercise~\ref{c6:ex:1}
first, since Exercise~\ref{c6:ex:2} backpropagates through the values it produces. Set~2 is a
program, and it isn't new: the file \file{lectures/L06/conv_layer/cpp/conv_demo.cpp} in the
companion repository already holds a skeleton of the struct you write, its helper functions, and
the \code{main()} that uses it, with a Makefile beside it. Get as far as you can with it here;
Chapter~\ref{c7:ch:layers} finishes it.

\textbf{How to check your work.} The worked solutions to Set~1 are in
Appendix~\ref{sol:ch:solutions}. Set~2 has no test suite; instead, the demo runs the numbers of
Set~1 and prints each result next to the value the hand calculation arrives at, so your
implementation is checked against your own arithmetic.

The same layer is implemented in the repository twice more, as references: in Python, in
\file{lectures/L06/conv_layer/python}, and in C, in \file{lectures/L06/conv_layer/c}, where the
layer is encapsulated behind an opaque struct whose accessor functions are only exposed through its
header. Try the exercise yourself before looking at them.
\end{aside}

% ------------------------------------------------------------------------------------------
\exerciseset{Hand-Training a CNN}
\label{c6:set:1}\label{c6:app:b}
We have the following convolutional neural network:
\begin{itemize}
  \item Input: \textbf{$4 \times 4$}.
  \item Conv layer:
  \begin{itemize}
    \item Kernel: \textbf{$2 \times 2$}.
    \item Stride = 1 (the kernel moves one pixel at a time).
    \item Zero-padding (we add zeros around the image so its size is preserved despite the
      extraction).
    \item Activation: \textbf{ReLU} (all negative values are replaced with zero).
    \item Learning rate: $\LR = 0.001$, i.e.\ 0.1\%.
  \end{itemize}
  \item Max pool:
  \begin{itemize}
    \item Pool size = 2.
    \item Stride = 2 (the pools are non-overlapping).
  \end{itemize}
  \item Flatten: $2 \times 2 \rightarrow 1 \times 4$.
\end{itemize}
The two exercises below work the whole training process through, step by step: feedforward
through all three layers, then backpropagation and optimization.

\calcexercise{Feedforward}
\label{c6:ex:1}
Assume we have the following input image, resembling the digit 0 made up of ones:
\[
  \begin{bmatrix}
    1 & 1 & 1 & 1 \\
    1 & 0 & 0 & 1 \\
    1 & 0 & 0 & 1 \\
    1 & 1 & 1 & 1
  \end{bmatrix}
\]
We use a kernel (a filter of weights) and a bias, with simple, fixed values chosen to make the
calculations easy to follow:
\[
  \text{kernel} = \begin{bmatrix} 0.2 & 0.4 \\ 0.6 & 0.8 \end{bmatrix}, \qquad
  \text{bias} = 0.5
\]

\task{a) Padding the input image}
The first step is \emph{zero-padding}: we add zeros around the image so its size is preserved even
after the kernel filter has been applied. This way, the conv layer's output has the same dimensions
as the input, $4 \times 4$ in this case.

The number of zeros $n$ we add on each side is computed as
\[
  n = \text{kernel\_size} / 2,
\]
where \code{kernel_size} is the size of the kernel, 2 in this example, i.e.\ one zero on each side.
Give the padded image, which is $6 \times 6$.

In general, sliding a $K \times K$ kernel with stride $s$ over an $N \times N$ image padded by $p$
on each side gives an output of size $\lfloor (N + 2p - K)/s \rfloor + 1$. For an \textbf{odd}
kernel that lands exactly on $N$, which is the case this padding rule is designed for. For the
\textbf{even} $2 \times 2$ kernel used here it comes out one larger, $5 \times 5$, because
\code{n = kernel_size / 2} pads symmetrically for a kernel that has no centre. The layer keeps only
the first $4 \times 4$ window positions, i.e.\ $i$ and $j$ both run from 0 to 3, which is what makes
the output match the input size. Work out that $4 \times 4$ below and you have exactly what the code
produces.

\task{b) Feedforward through the conv layer}
For each position in the image, compute the sum of the element-wise multiplication between the
kernel and the corresponding part of the padded image, plus the bias, then apply ReLU (all negative
values are replaced with zero):

\begin{console}
sum = a * k00 + b * k01 + c * k10 + d * k11 + bias
output = ReLU(sum)
\end{console}

\noindent where:
\begin{itemize}
  \item \textbf{a, b, c, d} are pixel values from the current $2 \times 2$ section of the image.
  \item \textbf{k} stands for kernel (the filter), e.g.\ \code{k00} means the value at row 0,
    column 0 of the kernel matrix, i.e.\ the top-left corner.
  \item \textbf{bias} is a constant value added to the sum.
\end{itemize}
Apply the kernel and bias across the entire image, and give the conv output ($4 \times 4$).

\task{c) The max pooling layer (feedforward)}
The max pooling layer looks for the largest value in each $2 \times 2$ pool (non-overlapping) of the
conv layer's output. The purpose is to downsample the image: we remove detail while keeping the most
important features, which reduces the amount of data and makes the model less sensitive to small
variations.

Split the conv layer's output into its four pools, and give the largest value of each, and the
position it was found at. Where a pool holds its maximum more than once, pass forward the first
instance. Then give the max pooling layer's output ($2 \times 2$).

\task{d) Flatten (feedforward)}
Flatten the max pooling layer's output into a vector, so it can be fed to the next layer.
\seesolution{c6:ex:1}

\calcexercise{Backpropagation and optimization}
\label{c6:ex:2}
Assume the dense layer sends back the following gradients:
\[
  [10, 20, 30, 40]
\]

\task{a) Flatten (backpropagation)}
Reshape the gradients from the dense layer back into a matrix.

\task{b) Max pool (backpropagation)}
Propagate the gradients back to the correct positions in the max pooling layer, i.e.\ to the spots
where the max values were located in each pool. If a pool holds its max value more than once, the
gradient is propagated back to the first position; the remaining positions get a gradient of 0.
Give the gradient for each pool, and then all four assembled into a single $4 \times 4$ matrix.

\task{c) Conv layer (backpropagation)}
Now compute the gradient for the bias and the kernel from the error propagated back.

One step is invisible here. Strictly, each incoming gradient is first multiplied by the activation
function's derivative at that position, $\sigma'(s)$, where $s$ is the pre-activation sum from
feedforward. Every sum in this example is positive and the activation is ReLU, so every derivative
is 1 and the multiplication changes nothing. It is not optional in code: the \code{Conv} class you
write in Chapter~\ref{c8:ch:conv} applies it explicitly, in the \code{backpropagate()} specified in
\secref{c8:sec:exercises}.

\textbf{Bias gradient:} the sum of all values in the gradient matrix.

\textbf{Kernel gradients:} for each kernel element, sum the product of the corresponding patch in
the padded input image and the gradient matrix. For each position $(i, j)$ in the gradient matrix,
the corresponding $2 \times 2$ patch is extracted from the padded input image $P$, and each kernel
element is multiplied by the corresponding value in the patch and summed over all positions:
\[
  \text{kernelGradient}[a][c] = \sum_{i,j} P[i+a][j+c] \cdot \text{grad}[i][j]
\]
Most gradients are zero, so only the positions where the gradient is non-zero need computing.

\textbf{Input gradients (padded input):} compute how the error propagates back to the input by, for
each position in the gradient matrix, ``spreading'' the kernel's weights multiplied by the gradient
value at the right spot into a new $6 \times 6$ matrix, and summing overlapping positions. Then
remove the padding, which leaves a $4 \times 4$ matrix matching the original image: the gradient
with respect to the input, which is passed further back through the network.

\task{d) Optimization}
Finally, update the kernel and bias using the learning rate $\LR$:

\begin{console}
kernel = kernel + LR * kernelGradient
bias   = bias   + LR * biasGradient
\end{console}

\noindent Give the updated kernel and the updated bias.
\seesolution{c6:ex:2}

% ------------------------------------------------------------------------------------------
\exerciseset{A Simple Conv Layer in C++}
\label{c6:set:2}\label{c6:app:c}

\codeexercise{The \code{ml::ConvLayer} struct}
\label{c6:ex:3}
A struct named \code{ml::ConvLayer} should be added to
\file{lectures/L06/conv_layer/cpp/conv_demo.cpp} to implement a simple conv layer. To keep things as
simple as possible, we implement a struct and skip get/set methods, deletion of copy and move
constructors, and so on.

You don't need to finish this now: get as far as you can with \code{feedforward()} and
\code{backpropagate()}; you'll finish the struct (and compile and run it for the first time) in
Chapter~\ref{c7:ch:layers}. See Exercise~\ref{c7:ex:1} for the wrap-up steps once you're done.

\code{ConvLayer} zero-pads its input (\code{pad = kernelSize / 2}) so the output is the
\textbf{same size as the input}, not smaller. This is why \code{outputGradients} in the example below
is a $4 \times 4$ matrix, matching the $4 \times 4$ \code{input}: your \code{feedforward()} and
\code{backpropagate()} need to pad and unpad internally to make that size match up. The struct
skeleton in \file{conv_demo.cpp} sketches two private helper methods for exactly that, along with
every member variable you need.

The general output size is $\lfloor (N + 2p - K)/s \rfloor + 1$, which lands exactly on $N$ for an
odd kernel. The $2 \times 2$ kernel used here is even, so the padded image admits $5 \times 5$ window
positions; loop \code{i} and \code{j} from 0 to \code{inputSize - 1} and the extra row and column
simply go unused. Exercise~\ref{c6:ex:1} works the same point through by hand.

The demo runs the \emph{same} numbers you worked through by hand in Exercises~\ref{c6:ex:1}
and~\ref{c6:ex:2}: the same input image, the same kernel and bias, and the same gradients coming
back from the max pooling layer. It prints each result next to the value the hand calculation
arrives at, so you can check your implementation against your own hand calculations rather than
guessing whether the output looks reasonable.

That's also why the kernel and bias are assigned right after the layer is created: the constructor
randomizes them, and the demo overwrites them so every run prints the same numbers.

Study the code in the \code{main()} function. Your implementation should be written so this code
works:

\begin{cppcode}
// Create a convolutional layer: 4x4 input, 2x2 kernel.
constexpr std::size_t inputSize{4U};
constexpr std::size_t kernelSize{2U};
ml::ConvLayer convLayer{inputSize, kernelSize};

// Use appendix B's kernel and bias instead of the randomized ones.
convLayer.kernel = Matrix2d{{0.2, 0.4}, {0.6, 0.8}};
convLayer.bias   = 0.5;

// Input image from appendix B, resembling the digit 0 made up of ones.
const Matrix2d input{{1, 1, 1, 1},
                     {1, 0, 0, 1},
                     {1, 0, 0, 1},
                     {1, 1, 1, 1}};

// Perform feedforward (convolution). Expect the conv output from appendix B.
convLayer.feedforward(input);

// Gradients the max pooling layer sends back in appendix B.
const Matrix2d outputGradients{{0, 10, 20, 0},
                               {0,  0,  0, 0},
                               {0,  0,  0, 0},
                               {0, 30,  0, 40}};

// Perform backpropagation. Expect the input gradients from appendix B.
convLayer.backpropagate(outputGradients);

// Perform optimization. Expect the updated kernel and bias from appendix B.
convLayer.optimize(0.001);
\end{cppcode}

\noindent The comments name appendix B of the lecture, which is Set~1 here. Each of
\code{feedforward()}, \code{backpropagate()} and \code{optimize()} returns a \code{bool}; the demo
checks it and aborts with a message on failure.
